\section{The \texttt{NCrystal\_process} McStas Component}
NCrystal\_process component for the Union framework.

\subsection*{Identification}
\begin{itemize}
  \item \textbf{Site:} 
  \item \textbf{Author:} NCrystal developers, converted to a Union component by Mads Bertelsen
  \item \textbf{Origin:} NCrystal Developers (European Spallation Source ERIC and DTU Nutech)
  \item \textbf{Date:} 20.08.15
\end{itemize}

\subsection*{Description}
\begin{lstlisting}
This process uses the NCrystal library as a Union process, see user documentation
for the NCrystal_sample.comp component for more information. The process only
uses the physics, as the Union components has a separate geometry system.
Absorption is also handled by Union, so any absorption output from NCrystal
is ignored.

Part of the Union components, a set of components that work together and thus
sperates geometry and physics within McStas.
The use of this component requires other components to be used.

1) One specifies a number of processes using process components like this one
2) These are gathered into material definitions using Union_make_material
3) Geometries are placed using Union_box / Union_cylinder, assigned a material
4) A Union_master component placed after all of the above

Only in step 4 will any simulation happen, and per default all geometries
defined before the master, but after the previous will be simulated here.

There is a dedicated manual available for the Union_components


Original header text for NCrystal_sample.comp:
McStas sample component for the NCrystal scattering library. Find more
information at <a href="https://github.com/mctools/ncrystal/wiki">the NCrystal
wiki</a>. In particular, browse the available datafiles at <a
href="https://github.com/mctools/ncrystal/wiki/Data-library">Data-library</a>
and read about format of the configuration string expected in the "cfg"
parameter at <a href="https://github.com/mctools/ncrystal/wiki/Using-NCrystal">Using-NCrystal</a>.

<p/>NCrystal is available under the <a href="http://www.apache.org/licenses/LICENSE-2.0">Apache 2.0 license</a>. Depending
on the configuration choices, optional NCrystal modules under different
licenses might be enabled - see <a
href="https://github.com/mctools/ncrystal/wiki/About">About</a> for more
details.

Algorithm:
Described elsewhere
\end{lstlisting}

\subsection*{Input parameters}
Parameters in \textbf{boldface} are required; the others are optional.

\begin{longtable}{p{0.22\textwidth}p{0.12\textwidth}p{0.46\textwidth}p{0.14\textwidth}}
\toprule
\textbf{Name} & \textbf{Unit} & \textbf{Description} & \textbf{Default} \\
\midrule
\endhead
cfg & str & NCrystal material configuration string (details \textless{}a href="https://github.com/mctools/ncrystal/wiki/Using-NCrystal"\textgreater{}on this page\textless{}/a\textgreater{}). & "" \\
packing\_factor & 1 & Material packing factor & 1 \\
interact\_fraction & 1 & How large a part of the scattering events should use this process 0-1 (sum of all processes in material = 1) & -1 \\
init & string & Name of Union\_init component (typically "init", default) & "init" \\
\bottomrule
\end{longtable}

\subsection*{Links}
\begin{itemize}
  \item \href{run:/home/willend/willend-McCode/mcstas-comps/union/NCrystal_process.comp}{Source code} for \texttt{NCrystal\_process.comp}.
\end{itemize}
\IfFileExists{NCrystal_process_static.tex}{\input{NCrystal_process_static.tex}}{}